\markdownRendererDocumentBegin
\markdownRendererSectionBegin
\markdownRendererHeadingOne{Discussion and conclusion}\markdownRendererInterblockSeparator
{}P12 reframes test-time scaling as a \markdownRendererStrongEmphasis{portfolio of computations}. “Think longer” is not the only adaptive action available to an intelligent system. It may be cheaper to parse, retrieve, compile, restructure or recover state so that less downstream search is required. Conversely, when state already exposes the relevant structure, additional preprocessing is wasteful and reasoning should receive the marginal budget.\markdownRendererInterblockSeparator
{}P12A demonstrates the construction but not the key discriminator: equal total budget was real while equal action capability was not. P12B repairs that controlled estimand and finds a positive two-signal effect across the registered equal-action panel.\markdownRendererInterblockSeparator
{}This leaves a concrete systems hypothesis for real agents: \markdownRendererStrongEmphasis{test-time scaling curves may be two-dimensional, with state-work and reasoning-work measured on a common receipt and with action capability held fixed across signal ablations.}\markdownRendererInterblockSeparator
{}Adaptive inference may need to decide not only \markdownRendererStrongEmphasis{how much} computation to spend but \markdownRendererStrongEmphasis{where} to spend it. P12 supplies the formulation, an exact failure analysis of its first empirical discriminator, and a positive prospectively frozen equal-action successor. Current authority is the bounded equal-action signal-complementarity result; the P12A comparison boundary remains historical, and matched real end-to-end validation is the next claim frontier.
\markdownRendererSectionEnd \markdownRendererDocumentEnd